%% ----------------------------------------------------------------------------
%%  §9.6  Integrazione per sostituzione
%% ----------------------------------------------------------------------------
\section{Integrazione per sostituzione}
\label{sec:sostituzione}

Il metodo di integrazione per decomposizione in somma si basa sulla regola di
derivazione della somma, quello per parti sulla regola di derivazione del prodotto. Il
metodo di integrazione per sostituzione si basa sulla regola di derivazione delle
funzioni composte.

\begin{teorema}[Formula di integrazione per sostituzione]\label{teo:sostituzione}
\index{integrazione!per sostituzione}
Sia $f$ una funzione continua in un intervallo $I$ e sia $g$ una funzione definita in un
intervallo $J$, a valori in $I$, derivabile con derivata continua. Allora
\begin{equation}\label{eq:sostituzione}
  \Bigl[\int f(x)\,dx\Bigr]_{x=g(t)}=\int f\bigl(g(t)\bigr)\,g'(t)\,dt ,
\end{equation}
nel senso che, se $F$ è una primitiva di $f$ in $I$, allora $F\bigl(g(t)\bigr)$ è una
primitiva di $f\bigl(g(t)\bigr)\,g'(t)$ in $J$.
\end{teorema}

\begin{schema}
\Ipotesi{$f$ continua nell'intervallo $I$; \ $g:J\to I$ derivabile con $g'$ continua.}
\Tesi{$F'=f \implica \bigl(F\circ g\bigr)'=(f\circ g)\,g'$; \ cioè la
\eqref{eq:sostituzione}.}
\end{schema}

\needspace{5\baselineskip}
\begin{strategia}
È la regola della catena letta al contrario: si deriva $F(g(t))$ e si trova
$F'(g(t))\,g'(t)=f(g(t))\,g'(t)$.
\end{strategia}

\begin{dimostrazione}
La funzione $f(g(t))\,g'(t)$ è continua in $J$, come prodotto della composta di funzioni
continue per una funzione continua, e quindi ha primitive. Se $F$ è una primitiva di $f$
in $I$, per il \cref{teo:derivata-composta}
\[
  \frac{d}{dt}F\bigl(g(t)\bigr)=F'\bigl(g(t)\bigr)\,g'(t)=f\bigl(g(t)\bigr)\,g'(t)
  \qquad\perogni t\in J ,
\]
cioè $F\circ g$ è una primitiva di $(f\circ g)\,g'$, come volevasi dimostrare.
\end{dimostrazione}

Il primo membro della \eqref{eq:sostituzione} significa: si calcola $\int f(x)\,dx=F(x)+c$
e poi si sostituisce $x=g(t)$, ottenendo $F(g(t))+c$. La formula non richiede che $g$ sia
invertibile; il risultato dell'integrazione, però, è espresso in funzione di $t$. Per
esprimerlo in funzione di $x$ occorre che $g$ sia invertibile (ad esempio strettamente
monotòna), e si effettua l'ulteriore sostituzione $t=g^{-1}(x)$. Come vedremo, per il
calcolo di un integrale \emph{definito} l'inversione non è necessaria.

La formula si ricorda facilmente con la notazione di Leibniz. Se $x=g(t)$, la quantità
$g'(t)\,dt$ si chiama \emph{differenziale} della funzione $g$ e si indica con $dx$:
\begin{equation}\label{eq:differenziale-sostituzione}
  dx=g'(t)\,dt
\end{equation}
(è il differenziale del \cref{sec:differenziabilita}, per una funzione di una variabile).
Nella \eqref{eq:sostituzione} la variabile $x$ si trasforma in $g(t)$ e il simbolo $dx$ si
trasforma secondo la \eqref{eq:differenziale-sostituzione}, come se fosse un prodotto. È
uno dei motivi del successo della notazione di Leibniz: le regole di calcolo sono
suggerite dai simboli stessi.

La formula si usa in due direzioni. Nella prima, già incontrata nella
\eqref{eq:composta-indefinito}, si riconosce nella funzione integranda la forma
$f(g(t))\,g'(t)$ e si pone $x=g(t)$ per ricondursi a $\int f(x)\,dx$. Nella seconda, si
sceglie una sostituzione $x=g(t)$ che renda più semplice la funzione integranda,
tipicamente eliminando una radice.

\begin{esempio}[Sostituzioni che eliminano una radice]\label{es:sostituzioni}
\begin{enumerate}[label=\textup{(\alph*)}]
  \item Nell'integrale $\int\frac{dx}{\sqrt x-3}$ è naturale porre $x=t^{2}$, con $t\ge0$
  e $t\neq3$, così che $dx=2t\,dt$ e $\sqrt x=t$:
  \[
  \begin{aligned}
    \int\frac{dx}{\sqrt x-3}&=\int\frac{2t}{t-3}\,dt=2\int\Bigl(1+\frac3{t-3}\Bigr)dt\\
    &=2t+6\log\abs{t-3}+c ,
  \end{aligned}
  \]
  e, tornando alla variabile $x$ con $t=\sqrt x$,
  \[
    \int\frac{dx}{\sqrt x-3}=2\sqrt x+6\log\abs{\sqrt x-3}+c .
  \]
  \item Con la sostituzione $2x-1=t^{2}$, cioè $x=\frac{t^{2}+1}2$ con $t\ge0$, si ha $dx=t\,dt$ e
  $x-\sqrt{2x-1}=\frac{(t-1)^{2}}2$, che si annulla per $t=1$, cioè $x=1$:
  \[
    \int\frac{x+\sqrt{2x-1}}{x-\sqrt{2x-1}}\,dx=\int\frac{(t^{2}+1)/2+t}{(t-1)^{2}/2}\,t\,dt
    =\int\frac{t^{3}+2t^{2}+t}{t^{2}-2t+1}\,dt .
  \]
  È l'integrale di una funzione razionale, con denominatore a radici coincidenti; la
  divisione e la \eqref{eq:fratti-semplici-doppia} danno
  \[
    \frac{t^{3}+2t^{2}+t}{(t-1)^{2}}=t+4+\frac8{t-1}+\frac4{(t-1)^{2}} ,
  \]
  e quindi l'integrale vale
  \[
    \frac{t^{2}}2+4t+8\log\abs{t-1}-\frac4{t-1}+c, \qquad\text{con } t=\sqrt{2x-1} .
  \]
  \item Con la sostituzione $x=\sen t$, $t\in[-\frac\pi2,\frac\pi2]$, si ha $dx=\cos t\,dt$ e
  $\sqrt{1-\sen^{2}t}=\cos t$, perché in questo intervallo $\cos t\ge0$; per la
  \eqref{eq:integrale-cos2},
  \begin{equation}\label{eq:integrale-semicerchio}
  \begin{split}
    \int\sqrt{1-x^{2}}\,dx&=\int\cos^{2}t\,dt=\frac12(t+\sen t\cos t)+c\\
    &=\frac12\Bigl(\arcsen x+x\sqrt{1-x^{2}}\Bigr)+c .
  \end{split}
  \end{equation}
  Nell'ultimo passaggio si è invertita la sostituzione: $t=\arcsen x$, $\cos t=\sqrt{1-x^{2}}$.
\end{enumerate}
\end{esempio}

\subsection{Sostituzione negli integrali definiti}

\begin{teorema}[Sostituzione negli integrali definiti]\label{teo:sostituzione-definita}
Sia $f$ continua in un intervallo $I$ e sia $g$ derivabile con derivata continua in
$[c,d]$, a valori in $I$, con
\begin{equation}\label{eq:estremi-sostituzione}
  g(c)=a, \qquad g(d)=b .
\end{equation}
Allora
\begin{equation}\label{eq:sostituzione-definita}
  \int_{a}^{b}f(x)\,dx=\int_{c}^{d}f\bigl(g(t)\bigr)\,g'(t)\,dt .
\end{equation}
\end{teorema}

\begin{schema}
\Ipotesi{$f$ continua in $I$; \ $g\in C^{1}([c,d])$, $g([c,d])\subseteq I$; \ $g(c)=a$,
$g(d)=b$.}
\Tesi{$\int_{a}^{b}f(x)\,dx=\int_{c}^{d}f(g(t))\,g'(t)\,dt$.}
\end{schema}

\begin{strategia}
Con una primitiva $F$ di $f$, entrambi i membri si calcolano con la formula fondamentale:
il primo vale $F(b)-F(a)$, il secondo, poiché $F\circ g$ è una primitiva della sua funzione
integranda, vale $F(g(d))-F(g(c))$. Le condizioni \eqref{eq:estremi-sostituzione} dicono che
i due numeri coincidono.
\end{strategia}

\begin{dimostrazione}
Sia $F$ una primitiva di $f$ in $I$, che esiste perché $f$ è continua. Per la formula
fondamentale, $\int_{a}^{b}f(x)\,dx=F(b)-F(a)$ (anche se $a>b$, per la
\eqref{eq:estremi-invertiti}). Per il \cref{teo:sostituzione}, $F\circ g$ è una primitiva
della funzione continua $(f\circ g)\,g'$ in $[c,d]$, e quindi, ancora per la formula
fondamentale e per le \eqref{eq:estremi-sostituzione},
\[
  \int_{c}^{d}f\bigl(g(t)\bigr)\,g'(t)\,dt=F\bigl(g(d)\bigr)-F\bigl(g(c)\bigr)=F(b)-F(a) . \qedhere
\]
\end{dimostrazione}

Si osservi che non si richiede che $g$ sia invertibile, né che sia monotòna: conta
soltanto che $g$ porti gli estremi $c$, $d$ negli estremi $a$, $b$. Con $g(t)=\alpha t$ o
$g(t)=t+\tau$ si ritrovano, per le funzioni continue, le formule di cambiamento di scala e
di traslazione dell'\cref{es:scala-traslazione}; con $g(t)=-t$ quella di riflessione
dell'\cref{es:riflessione}.

\begin{esempio}[L'area del semicerchio]\label{es:semicerchio}
Calcoliamo $\int_{-1}^{1}\sqrt{1-x^{2}}\,dx$ con la sostituzione $x=\sen t$. Poiché
$\sen(-\frac\pi2)=-1$ e $\sen\frac\pi2=1$, per la \eqref{eq:sostituzione-definita} e la
\eqref{eq:integrale-cos2},
\begin{equation}\label{eq:area-semicerchio}
\begin{split}
  \int_{-1}^{1}\sqrt{1-x^{2}}\,dx&=\int_{-\pi/2}^{\pi/2}\cos^{2}t\,dt\\
  &=\Bigl[\frac12(t+\sen t\cos t)\Bigr]_{-\pi/2}^{\pi/2}=\frac\pi2 .
\end{split}
\end{equation}
È la verifica promessa nell'\cref{es:portanza-ellittica-esatta}: l'area del semicerchio
unitario definita mediante l'integrale è $\frac\pi2$, in accordo con la definizione di area
del cerchio mediante i poligoni inscritti (\cref{teo:area-cerchio}). Di conseguenza la
portanza totale della distribuzione ellittica è davvero $L=\frac\pi4\,l_{0}b$.
\end{esempio}

\begin{python}[Verifica delle sostituzioni]
import sympy as sp

x = sp.symbols("x", real=True)
F = (sp.asin(x) + x * sp.sqrt(1 - x**2)) / 2       # formula trovata
d = sp.simplify(sp.diff(F, x) - sp.sqrt(1 - x**2))
print("F' - sqrt(1-x^2) =", d)
print("int_-1^1 sqrt(1-x^2) dx =",
      sp.integrate(sp.sqrt(1 - x**2), (x, -1, 1)))
t = sp.sqrt(2*x - 1)                                # esempio (b)
G = t**2 / 2 + 4*t + 8 * sp.log(sp.Abs(t - 1)) - 4 / (t - 1)
f = (x + sp.sqrt(2*x - 1)) / (x - sp.sqrt(2*x - 1))
for x0 in (sp.Rational(3, 4), 3, 7):
    d = sp.simplify(sp.diff(G, x).subs(x, x0) - f.subs(x, x0))
    print(f"x = {x0}:  G'(x) - f(x) = {sp.nsimplify(d)}")
\end{python}
\begin{risultato}
F' - sqrt(1-x^2) = 0
int_-1^1 sqrt(1-x^2) dx = pi/2
x = 3/4:  G'(x) - f(x) = 0
x = 3:  G'(x) - f(x) = 0
x = 7:  G'(x) - f(x) = 0
\end{risultato}

\begin{notastorica}[Il differenziale di Leibniz]
Per Leibniz $dx$ era una quantità «infinitamente piccola», e l'integrale una somma di
infiniti termini $f(x)\,dx$: in questa visione la regola di sostituzione è quasi ovvia,
perché un rettangolino di base $dx$ si trasforma in uno di base $g'(t)\,dt$. La teoria dei
limiti di Cauchy e Weierstrass eliminò gli infinitesimi attuali, e il simbolo $dx$
sopravvisse come pura notazione, giustificata dai teoremi di questo paragrafo. Soltanto
nel 1961 Abraham Robinson mostrò, con l'\emph{analisi non standard} (esposta poi nel libro
omonimo del 1966), che gli infinitesimi di Leibniz si possono rendere rigorosi, ampliando
i numeri reali con numeri infinitamente piccoli e infinitamente grandi.
\end{notastorica}

\subsection{Approfondimento: alcune sostituzioni standard}
\label{subsec:sostituzioni-standard}

Per alcune classi di funzioni esistono sostituzioni che riconducono \emph{sempre}
l'integrale a quello di una funzione razionale (\cref{sec:razionali}). Ne presentiamo
due, seguendo \cite{giusti}. Indichiamo con $R(u,v)$ una \emph{funzione razionale di due
variabili}, quoziente di due polinomi in $u$ e $v$.

\paragraph{Funzioni razionali di $\sen x$ e $\cos x$.} Si pone $t=\tg\frac x2$.

\begin{proposizione}[La sostituzione $t=\tg\frac x2$]\label{prop:sostituzione-tg-mezzi}
Per $x\in(-\pi,\pi)$, posto $t=\tg\frac x2$, cioè $x=2\operatorname{arctg}t$, si ha
\begin{equation}\label{eq:formule-parametriche}
  \sen x=\frac{2t}{1+t^{2}}, \qquad \cos x=\frac{1-t^{2}}{1+t^{2}}, \qquad
  dx=\frac{2\,dt}{1+t^{2}} .
\end{equation}
Di conseguenza, se $f(x)=R(\sen x,\cos x)$ è continua in un intervallo $J\subseteq(-\pi,\pi)$,
$K$ è l'intervallo descritto da $\tg\frac x2$ per $x\in J$, e $\Phi$ è una primitiva in $K$
della funzione razionale
\begin{equation}\label{eq:integranda-razionale-t}
  r(t)=R\Bigl(\frac{2t}{1+t^{2}},\frac{1-t^{2}}{1+t^{2}}\Bigr)\,\frac2{1+t^{2}} ,
\end{equation}
allora $\Phi\bigl(\tg\frac x2\bigr)$ è una primitiva di $f$ in $J$.
\end{proposizione}

\begin{schema}
\Ipotesi{$J\subseteq(-\pi,\pi)$ intervallo; \ $f=R(\sen x,\cos x)$ continua in $J$; \ $\Phi'=r$.}
\Tesi{le \eqref{eq:formule-parametriche}; \ $\frac{d}{dx}\Phi\bigl(\tg\frac x2\bigr)=f(x)$
per ogni $x\in J$.}
\end{schema}

\begin{strategia}
Si usano le formule di duplicazione nell'angolo $\frac x2$, divise per
$\cos^{2}\frac x2+\sen^{2}\frac x2=1$, così da far comparire soltanto la tangente. Poi si
deriva $\Phi(\tg\frac x2)$ con la regola della catena: il fattore $\frac2{1+t^{2}}$ di $r$ si
semplifica con la derivata di $\tg\frac x2$.
\end{strategia}

\begin{dimostrazione}
Se $x\in(-\pi,\pi)$, allora $\frac x2\in(-\frac\pi2,\frac\pi2)$ e $\cos\frac x2>0$, quindi
$t=\tg\frac x2$ è definita. Per le formule di duplicazione \eqref{eq:duplicazione},
dividendo numeratore e denominatore per $\cos^{2}\frac x2$,
\[
\begin{gathered}
  \sen x=\frac{2\sen\frac x2\cos\frac x2}{\cos^{2}\frac x2+\sen^{2}\frac x2}
  =\frac{2t}{1+t^{2}},\\
  \cos x=\frac{\cos^{2}\frac x2-\sen^{2}\frac x2}{\cos^{2}\frac x2+\sen^{2}\frac x2}
  =\frac{1-t^{2}}{1+t^{2}} .
\end{gathered}
\]
La funzione $h(x)=\tg\frac x2$ è strettamente crescente in $(-\pi,\pi)$, con inversa
$2\operatorname{arctg}t$ (la cui derivata $\frac2{1+t^{2}}$ è il fattore $dx$ della
\eqref{eq:formule-parametriche}) e derivata $h'(x)=\frac12\bigl(1+\tg^{2}\frac x2\bigr)
=\frac{1+t^{2}}2$. La funzione $r(t)=f(2\operatorname{arctg}t)\frac2{1+t^{2}}$ è continua in $K$
e, per il \cref{teo:derivata-composta}, per ogni $x\in J$,
\[
  \frac{d}{dx}\Phi\bigl(h(x)\bigr)=r\bigl(h(x)\bigr)\,h'(x)
  =f(x)\,\frac2{1+t^{2}}\cdot\frac{1+t^{2}}2=f(x) . \qedhere
\]
\end{dimostrazione}

\begin{esempio}[Una primitiva che vale soltanto in un intervallo]\label{es:due-piu-coseno}
Per $f(x)=\frac1{2+\cos x}$, che è continua in $\R$, la
\eqref{eq:integranda-razionale-t} dà
$r(t)=\frac{1+t^{2}}{3+t^{2}}\cdot\frac2{1+t^{2}}=\frac2{3+t^{2}}$, e per la
\eqref{eq:argomento-lineare}
\[
  \int\frac{dx}{2+\cos x}=\frac2{\sqrt3}\operatorname{arctg}\Bigl(\frac1{\sqrt3}\tg\frac x2\Bigr)+c
  \qquad\text{in }(-\pi,\pi) .
\]
La primitiva trovata, $G$, \emph{non} vale in tutto $\R$: per $x\to\pi^{-}$ tende a
$\frac\pi{\sqrt3}$, per $x\to\pi^{+}$ (in $(\pi,3\pi)$ la stessa formula è ancora una
primitiva) tende a $-\frac\pi{\sqrt3}$. Il calcolo $\bigl[G(x)\bigr]_{0}^{2\pi}=0$ è
quindi sbagliato, come mostra il segno positivo della funzione integranda; quello
corretto spezza l'intervallo:
la funzione integrale $F(x)=\int_{0}^{x}f$ è continua in $[0,\pi]$ e, essendo come $G$ una
primitiva di $f$ in $[0,\pi)$ nulla in $0$, coincide con $G$ in $[0,\pi)$
(\cref{teo:caratterizzazione-primitive}); quindi
$\int_{0}^{\pi}f=F(\pi)=\lim_{x\to\pi^{-}}G(x)=\frac\pi{\sqrt3}$. Poiché
$\cos(2\pi-x)=\cos x$, la sostituzione $x=2\pi-u$ dà $\int_{\pi}^{2\pi}f=\int_{0}^{\pi}f$, e
in conclusione $\int_{0}^{2\pi}\frac{dx}{2+\cos x}=\frac{2\pi}{\sqrt3}$.
\end{esempio}

\paragraph{Funzioni razionali di $x$ e $\sqrt{x^{2}+a}$.} Qui l'idea, dovuta a Eulero, è
di eliminare la radice imponendo $\sqrt{x^{2}+a}=t-x$: elevando al quadrato il termine
$x^{2}$ si cancella, e resta un'equazione di primo grado in $x$.

\begin{proposizione}[La sostituzione di Eulero]\label{prop:sostituzione-eulero}
Siano $a\neq0$ e $J$ un intervallo in cui $x^{2}+a>0$. La funzione
$t=x+\sqrt{x^{2}+a}$ non si annulla ed è strettamente monotòna in $J$, e
\begin{equation}\label{eq:sostituzione-eulero}
  x=\frac{t^{2}-a}{2t}, \qquad \sqrt{x^{2}+a}=\frac{t^{2}+a}{2t}, \qquad
  dx=\frac{t^{2}+a}{2t^{2}}\,dt ;
\end{equation}
la sostituzione trasforma quindi $\int R\bigl(x,\sqrt{x^{2}+a}\bigr)\,dx$ nell'integrale
di una funzione razionale di $t$. In particolare,
\begin{equation}\label{eq:integrale-eulero}
  \int\frac{dx}{\sqrt{x^{2}+a}}=\log\Bigl|x+\sqrt{x^{2}+a}\Bigr|+c
  \qquad (x\in J) .
\end{equation}
\end{proposizione}

\begin{schema}
\Ipotesi{$a\neq0$; \ $J$ intervallo con $x^{2}+a>0$; \ $t(x)=x+\sqrt{x^{2}+a}$.}
\Tesi{$t\neq0$ e $t$ strettamente monotòna in $J$; \ le \eqref{eq:sostituzione-eulero}; \
la primitiva \eqref{eq:integrale-eulero}.}
\end{schema}

\begin{strategia}
Si risolve l'equazione di primo grado descritta sopra. La derivata di $t(x)$ risulta
proporzionale a $t$ stessa: ciò dà la monotonia, che serve per tornare alla variabile
$x$, e mostra che $\frac1{\sqrt{x^{2}+a}}$ è del tipo $\frac{\varphi'}\varphi$.
\end{strategia}

\begin{dimostrazione}
Se fosse $t=0$, cioè $\sqrt{x^{2}+a}=-x$, elevando al quadrato si avrebbe $a=0$: dunque
$t\neq0$ in $J$. Da $\sqrt{x^{2}+a}=t-x$ si ricava $x^{2}+a=t^{2}-2tx+x^{2}$, cioè
$x=\frac{t^{2}-a}{2t}$, e quindi $\sqrt{x^{2}+a}=t-x=\frac{t^{2}+a}{2t}$. La funzione $t(x)$ è
derivabile in $J$, con
\[
  t'(x)=1+\frac x{\sqrt{x^{2}+a}}=\frac{t(x)}{\sqrt{x^{2}+a}} .
\]
Poiché $t$ è continua e non si annulla nell'intervallo $J$, per il teorema degli zeri
(\cref{teo:zeri}) ha segno costante, e così $t'$: per il \cref{cor:derivata-positiva}
(applicato a $t$ o a $-t$), $t$ è strettamente monotòna. La sua inversa
$x(t)=\frac{t^{2}-a}{2t}$ ha derivata $\frac12+\frac a{2t^{2}}=\frac{t^{2}+a}{2t^{2}}$, che è il
fattore $dx$ della \eqref{eq:sostituzione-eulero}; per il \cref{teo:sostituzione}, con
$g(t)=\frac{t^{2}-a}{2t}$, la funzione integranda diventa
$R\bigl(\frac{t^{2}-a}{2t},\frac{t^{2}+a}{2t}\bigr)\frac{t^{2}+a}{2t^{2}}$, razionale in $t$.
Infine, per la seconda delle \eqref{eq:integrali-composti} con $\varphi=t(x)$,
\[
  \int\frac{dx}{\sqrt{x^{2}+a}}=\int\frac{t'(x)}{t(x)}\,dx=\log\abs{t(x)}+c . \qedhere
\]
\end{dimostrazione}

\begin{esempio}[L'integrale di $\sqrt{x^{2}+1}$]\label{es:radice-x2-piu-1}
Con $a=1$ si ha $t=x+\sqrt{x^{2}+1}>0$ e, per le \eqref{eq:sostituzione-eulero},
\[
  \sqrt{x^{2}+1}\,dx=\frac{(t^{2}+1)^{2}}{4t^{3}}\,dt
  =\Bigl(\frac t4+\frac1{2t}+\frac1{4t^{3}}\Bigr)dt ,
\]
la cui primitiva è $\frac{t^{2}}8+\frac12\log t-\frac1{8t^{2}}$. Per tornare alla variabile
$x$ conviene osservare che $t-\frac1t=\frac{t^{2}-1}t=2x$ e $t+\frac1t=2\sqrt{x^{2}+1}$, così che
$\frac{t^{2}}8-\frac1{8t^{2}}=\frac18\bigl(t-\frac1t\bigr)\bigl(t+\frac1t\bigr)=\frac12x\sqrt{x^{2}+1}$:
\[
  \int\sqrt{x^{2}+1}\,dx=\frac12\Bigl(x\sqrt{x^{2}+1}+\log\bigl(x+\sqrt{x^{2}+1}\bigr)\Bigr)+c ,
\]
analoga alla \eqref{eq:integrale-semicerchio}, con il logaritmo al posto
dell'arcoseno.
\end{esempio}

Un radicando di secondo grado qualsiasi, $\alpha x^{2}+\beta x+\gamma$, si riconduce ai
casi precedenti completando il quadrato, come nella \cref{prop:razionali-secondo-grado}.
Se $\alpha>0$, con $u=\sqrt\alpha\,\bigl(x+\frac\beta{2\alpha}\bigr)$ esso diventa $u^{2}+a$, e si
applica la sostituzione di Eulero; se $\alpha<0$, con $u=\sqrt{-\alpha}\,\bigl(x+\frac\beta{2\alpha}\bigr)$
diventa $k^{2}-u^{2}$, e la sostituzione $u=k\sen\theta$, come nell'\cref{es:sostituzioni}(c),
porta a una funzione razionale di $\sen\theta$ e $\cos\theta$, cioè al primo caso. Sostituzioni di entrambi i tipi compaiono
già nelle \emph{Institutiones calculi integralis} di Eulero (1768); quella con
$\tg\frac x2$ è oggi spesso attribuita, senza riscontri documentari, a Weierstrass.

